\documentclass[11pt]{article}
\usepackage[margin=1in]{geometry}
\usepackage[T1]{fontenc}
\usepackage{lmodern,microtype,hyperref,enumitem,amsmath,xcolor}
\hypersetup{colorlinks=true,urlcolor=blue,citecolor=blue,linkcolor=blue}
\title{Q01: How Can Next-Token Prediction Lead to Problem Solving, Coding, and Instruction Following?}
\author{DLS 2026 --- Draft module}
\date{October 2026}
\begin{document}
\maketitle

\section*{Learning objectives}
After this module, a student should be able to (i) state the next-token objective precisely, (ii) explain why the objective can reward modeling structure in text and code, (iii) distinguish observed capability from a settled account of mechanism, and (iv) propose an evaluation that separates memorization, pattern completion, and generalization.

\section{The narrow training objective}
An autoregressive language model assigns a probability to each next token conditional on preceding tokens.  Given a token sequence $x_{1:T}$, training commonly minimizes
\[
  \mathcal{L}(\theta)=-\sum_{t=1}^{T-1}\log p_\theta(x_{t+1}\mid x_{\leq t}).
\]
The objective does \emph{not} explicitly instruct the model to reason, write programs, follow an instruction, or represent a world model. It asks for a probability distribution that predicts continuations in the training distribution. A Transformer provides a flexible way to compute that distribution from the preceding sequence using attention and feed-forward layers \cite{vaswani2017attention}.

This apparent mismatch is the source of the question. The useful first answer is not that next-token prediction secretly contains a reasoning rule. Rather, many continuations are hard to predict without tracking regularities that humans describe as syntax, facts, goals, causal relations, conventions, and intermediate steps. If a document says ``the function returns the sum of the list,'' continuations in code, tests, and prose constrain which statements and tokens are likely next. A model that represents useful latent regularities can reduce prediction loss across many such contexts.

\section{Why broad capabilities can be useful for prediction}
Three routes help explain why the single objective can support varied behavior.

\paragraph{Compression of repeated structure.} Natural language and code contain repeated patterns at several scales: spelling, grammar, APIs, algorithms, argument structures, and task formats. A representation that captures a reusable pattern can improve predictions in many places. This is an explanation of \emph{why} representations may be useful; it is not proof that a particular internal feature carries a human-readable concept.

\paragraph{Learning from the prompt.} At inference time the prompt is part of the model's context. Demonstrations, constraints, and a question can therefore alter the next-token distribution without changing model weights. Large language models showed substantial few-shot behavior under this setup \cite{brown2020language}. In a controlled family of small Transformers, induction heads were associated with a sharp increase in an in-context learning behavior \cite{olsson2022induction}. This is a concrete mechanistic clue, not a universal explanation of all in-context learning.

\paragraph{Training data contains tasks and solutions.} Text and code corpora include explanations, debugging conversations, specifications, proofs, examples, and instruction-like exchanges. The model is rewarded for predicting the continuation of those patterns. Post-training can further alter which helpful, safe, or instruction-following continuation is preferred; Q01 concerns the pre-training objective and should not collapse this distinction.

\section{What this does \emph{not} establish}
A fluent continuation is not itself evidence of reliable reasoning. The same model can produce a correct multi-step derivation, a plausible but false derivation, or a memorized fragment. Evaluations must therefore test held-out distributions, vary surface form, require verifiable intermediate artifacts where appropriate, and report failure modes.

Likewise, ``emergence'' is a description of an observed performance pattern, not a complete causal explanation. Some capabilities can look discontinuous when a continuous score is passed through a thresholded metric; this criticism is developed by Schaeffer et al. \cite{schaeffer2023mirage}. Other work documents sharp-looking changes on selected tasks as scale increases \cite{wei2022emergent}. Students should ask: What metric was plotted? Is the effect robust to a different metric, model family, prompt, and contamination control?

\section{A worked miniature: copying a latent rule}
Consider the sequence \texttt{A B A B A}. Predicting its next token is easy if a model notices the repeating two-token pattern; predicting it by a fixed unigram frequency is not. The accompanying program implements two deliberately small predictors: a unigram baseline and an ``induction-like'' copy heuristic that finds an earlier occurrence of the current context and copies the token that followed it. The heuristic is not a Transformer and does not prove that real models work this way. It makes one key idea inspectable: a context-dependent rule can beat a context-free rule on a task whose continuation depends on latent structure.

\section{Takeaway}
Next-token prediction is a narrow objective applied to a broad, structured corpus. It can reward internal representations and context-sensitive algorithms that help predict language, code, and task demonstrations. That gives a plausible route to broad capabilities; it does not settle whether a response reflects robust reasoning, how a particular model implements it, or whether the behavior will generalize. Those are empirical questions requiring careful evaluation.

\bibliographystyle{plain}
\bibliography{bibliography}
\end{document}
